%% ----------------------------------------------------------------
%% Chapter 1 -- Introduction
%% ----------------------------------------------------------------
\chapter{Introduction}
\label{chap:intro}

The integration of both the sensor functionalities with communication capabilities in wireless networks is referred to as Integrated Sensing and Communication (ISAC). In that sense, as sensing functionality is increasingly absorbed by network infrastructure, it will provide distributed physical layer awareness. Distributed physical-layer awareness will be particularly important in GPS-denied indoor environments because satellite signals are degraded due to attenuation caused by building structures. Nonetheless, there exist significant limitations in terms of physics when deploying single modality ISAC systems.

For example, while millimeter wave frequency modulated continuous wave (FMCW) radar (e.g., 24~GHz) has a fine range resolution (sub-foot), can estimate the velocity of targets via micro-Doppler effects, and has fine angular resolution (within its Field of View; $\pm$60 degrees), there is no identity associated with these types of systems. Radar returns cannot differentiate between a cooperative asset, a person, an object of clutter, or machinery. Additionally, because millimeter wave signals do not diffract well, when a Line-Of-Sight Occlusion occurs all target information is lost, and Multipath Reflected Signals create ``ghost'' detections.

On the other hand, IEEE 802.11mc Fine Time Measurement (FTM) at frequencies of 2.4~GHz and/or 5~GHz provide distance to a target (radial range) from time stamps associated with round trip times (RTTs) of packets sent from a mobile station to an access point based on the MAC address of the mobile station~\cite{aggarwal_wifi}. Therefore, this method inherently identifies the mobile station making the measurement and allows for non-line-of-sight detection due to better penetration of walls compared to millimeter wave. However, fine time measurement has spatial variability ranging from $\pm$1.5 to $\pm$4.8~m in the near field. FTM does not provide any angle information about the location of the mobile station relative to the access point, and updates occur every 1 to 2 seconds.

Fusion of these two complementary technologies leads to what we will call the synchronization gap. Millimeter wave radar produces data in the form of UART streamed frames from 10--24~Hz, while FTM functions at the level of the MAC layer as bursts, or processes at approximately 1~Hz. Kalman Filters rely upon synchronous updates and therefore inherently trust stale measurements when there are no new measurements available during intervals between updates (a condition referred to as ``freshness blindness''). As a result, the Kalman Filter exhibits Trajectory Divergence due to the Staircase Effect; that is, it locks onto old Wi-Fi readings during maneuvering by the Target.

This Work addresses the Fusion Problem for ISAC through the use of a Dual-Tier Framework running on Edge Hardware: An ESP32-S3 Dual Core Microprocessor running side-by-side with a HLK-LD2450 24~GHz FMCW Radar. The cooperative target consists of an active tag that includes both an ESP32 FTM responder as well as a trihedral corner reflector to increase Radar Cross Section.

We compare two different methods for fusing the radar and Wi-Fi sensors. Iteration~1 is dynamic covariance intersection (DCI). DCI is an information aware method that uses the angular Vector of the Radar to project the one dimensional Wi-Fi Range to geographically match it, performs First Order Hold Velocity Extrapolation to remove the Temporal Staircase effect, and dynamically increases the Stale Sensor Covariance using Age of Information penalty values~\cite{noroozi_aoi}. During radar occlusions, DCI allows trust to transition proportionally toward Wi-Fi without requiring any pre-programmed switching logic.

Iteration~2 was the Edge MLP. It is a multilayer perceptron model (11.6~KB float32) implemented using TensorFlow Lite Micro on the ESP32-S3 inside the SRAM. The input to the network is a $5 \times 5$ sliding window of forward filled feature values as well as binary freshness flags, learning velocity, acceleration and multi path distortion characteristics all derived from offline data. To prevent regression degradation due to INT8 quantization, we retained full Float32 precision. As a result, the MLP can achieve sub-10~cm accuracy tracking errors. Validation utilized both a MATLAB/Simulink-based digital twin model, modeling continuous time physics of the targets and discrete, multi-rate output from sensors, as well as a Python based test suite to simulate various occlusions, field of view escape and packet loss.
